\section{The spectrum of four-interface transmission}
\label{sp:section}
\label{sec:spectrum}

We now estimate the limiting transmission probabilities
$T^\pm(p,x,y)$ of Section~\ref{sec:annuli} when the modulus $p$ is large.
There are two different conclusions.  The leading term separates into an
inner factor and an outer factor.  The signed difference $T^+-T^-$ has
additional decay, and averaging the rotation of one end improves that
decay beyond exponent $2$.  This last improvement will make the response
to a fair Poisson insertion integrable over the plane.

\subsection{The modulus equation}

The space of four cyclically ordered, labeled angles is
\[
 \mathcal C=
 \{w\in\R^4:w_1<w_2<w_3<w_4<w_1+2\pi\}
       /(2\pi\mathbb Z\mathbf 1).
\]
Thus we quotient by simultaneous full turns, not by cyclic relabeling.
We use the product chamber $\mathcal X=\mathcal C\times\mathcal C$ for
$(x,y)$.  All statements about distances or derivatives on this chamber
use local angle lifts.  Constants are uniform as angles approach a
collision; moduli are bounded below by a fixed positive constant.

\begin{proposition}[Modulus equation]\label{sp:annular-pde}
The functions $T^\pm$ are smooth on $(0,\infty)\times\mathcal X$ and
satisfy
\begin{equation}\label{sp:pde}
 \partial_pT=(L+B_p\cdot\nabla)T,
 \qquad
 L=\frac18\sum_{w=x,y}
 \left(3\sum_{i=1}^4\partial_{w_iw_i}
 +\sum_{i=1}^4\sum_{j\ne i}
       \cot\frac{w_i-w_j}{2}\,\partial_{w_i}\right).
\end{equation}
For every $p_0>0$ and every fixed finite derivative order $k$,
the vector field $B_p$ extends smoothly across angle collisions and
\[
 \|B_p\|_{C^k}\le C_{k,p_0}e^{-p}\qquad(p\ge p_0).
\]
It is equivariant under permutations within either tuple, common
rotations, and cyclic relabelings.  In particular, its normal component
at every collision face is zero.
\end{proposition}

\begin{proof}
We first derive a local equation with just one switch active.  Write the
annulus as the cylinder $0<\operatorname{Im}z<p$, with real period
$2\pi$, and take the active switch on its lower boundary.  Parametrize
its exploration by the decrease of the modulus.  In the annulus Loewner
normalization with zero regular constant at the driving pole, the
angular kernel is
\[
 H_p(u)=\lim_{N\to\infty}\sum_{k=-N}^N
                     \cot\frac{u-2ikp}{2}.
\]
The pole has residue $2$, and the imaginary part on the other boundary
is $-1$.  The locality theorem for annulus $\mathrm{SLE}_6$
\cite{Zhan} identifies the local driving angle with Brownian motion
of variance parameter $6$.  Only locality before leaving a small
simply connected neighborhood of the active switch is used here.

For completeness, the absence of an additional driving drift depends
on the normalization.  In a chordal chart the pole part is
$2/(z-U_t)$.  If $F_t$ is a changing local conformal coordinate, comparison
of pole residues gives the new clock $F_t'(U_t)^2\,dt$.  Comparison of
the regular terms, both normalized to have zero constant, gives
$\partial_tF_t(U_t)=-3F_t''(U_t)$.  This cancels the It\^o correction
$\kappa F_t''(U_t)/2$ precisely for $\kappa=6$.
No condition on a remote target in the annulus enters this calculation.

Other switches on the active boundary have velocity $H_p(u)$, at
angular difference $u$ from the active one.  Switches on the opposite
boundary have angular velocity $\operatorname{Re}H_p(u+ip)$.
Pairing the positive and negative terms of the series gives, uniformly
with any fixed number of derivatives,
\begin{equation}\label{sp:kernel-errors}
 H_p(u)-\cot(u/2)=O(e^{-2p}),
 \qquad \operatorname{Re}H_p(u+ip)=O(e^{-p}).
\end{equation}
The remainders are smooth periodic functions of $u$, including at
$u=0$.  Thus, when $x_i$ is the active switch, its equation is
\begin{equation}\label{sp:active-equation}
 \begin{aligned}
 0={}&-\partial_pT+3\partial_{x_ix_i}T
       +\sum_{j\ne i}H_p(x_j-x_i)\,\partial_{x_j}T\\
    &+\sum_{j=1}^4\operatorname{Re}H_p(y_j-x_i+ip)\,
                                      \partial_{y_j}T.
 \end{aligned}
\end{equation}
For an active switch $y_i$, the equation is obtained by exchanging
$x$ and $y$.  Indeed the cylinder change $z\mapsto-z+ip$ exchanges
the boundaries and reverses both angle coordinates; the angular
kernels are odd, so transforming back preserves the displayed signs.
Averaging these eight equations gives \eqref{sp:pde}.
The series estimates give the asserted bounds on $B_p$.  The averaged
expression treats the labels within each tuple symmetrically; therefore
it is permutation equivariant.  At a collision its two corresponding
components agree, proving tangency.  The other symmetries follow in
the same way from the kernels.

We justify the Markov identity underlying these equations before
assuming differentiability.  On a round annulus choose an arbitrarily
small crosscut around the active switch, separated from all other
switches and from the other boundary component.  Approximate the smooth
bank by chord-free graph boundaries and complete the neighborhood to a
Jordan patch, putting its second color switch remotely.  Up to the
pre-query stop at the crosscut, the exploration in this patch is exactly
the exploration in the annulus: no color beyond the crosscut is read.
The stopped-exploration results of \cite[Lemmas~3.6--3.7]{V} give convergence
of the path together with its remaining domains.  Crosscuts may be
chosen as the semicircle crosscuts in the stopped-target construction;
the probability of an endpoint stop vanishes when endpoint
neighborhoods shrink.

The unaltered part, obtained by cutting off the fixed crosscut
neighborhood, is a Jordan annular back with unchanged boundary
traversal.  For the remaining face containing this back, the stopped-target-face
and remaining-domain statements of \cite[Lemmas~3.6--3.7]{V}
give finitely many Jordan side pieces attached to the back along
open subintervals of that crosscut.  Components already disconnected
from this face are discarded; no assertion of finiteness is being
made about all components swallowed by the exploration.
We check that its limit is still doubly connected.  In the closed
crosscut neighborhood, remove the side-lobe interiors and their open
attaching intervals.  The remaining closed set is connected to the
original outer complement.  Indeed join any of its points to the old
outer bank by a path in that neighborhood whose interior avoids the
crosscut.  If this path enters a lobe, replace its portion from first
entry to last exit by the corresponding subarc of the lobe boundary
off its open attachment.  The entry and exit lie on this retained
bank, not on the crosscut.  This bank meets the crosscut only at its
own endpoints and enters no other lobe interior, since the Jordan
interiors are disjoint.  Thus each replacement avoids all lobe
interiors and open attachments; finitely many replacements give a
path in the complement.  The unchanged inner hole is separated from
this connected outer complement by the persisting annular back.
These are precisely the two nondegenerate complementary components.
Lemma~\ref{an:back-lobes} gives convergence of annular uniformizations
on their closures, with the prime-end traversals and marks in order.
The hole and a separating open bulk persist throughout this local
exploration, so the moduli remain nondegenerate.  Thus the residual
annuli satisfy the hypotheses of Proposition~\ref{an:limit}.  The stopping
edge becomes the replacement switch, the other marks are unchanged,
and the colors in the remaining face are fresh.  Transmission topology
and bounded convergence therefore give the local mean-value identity
for $T^\pm$ at each of these hull stops.

This identity first implies the active equations in the viscosity
sense.  Work modulo simultaneous rotation and let a smooth test
function touch $T$ from above or below.  On a sufficiently small
crosscut stop the modulus and image marks remain in the touching
neighborhood.  Indeed the hull capacity tends to zero, and annular
maps normalized at a remote mark tend to the identity on the other
compact bank intervals, by reflection.  The image of the tip is
squeezed between neighboring unaffected bank points.  Applying It\^o's
formula to the touching function contradicts a strict inequality with
the wrong sign in the active equation, because the stop takes positive
time.  The same touching function can be used for all eight active
equations, so their average also holds in viscosity form.  On every
compact subset of the open chamber the averaged equation is a smooth,
uniformly parabolic linear equation, also after the rotation quotient.
Interior parabolic regularity then gives smoothness and the classical
equation.  Simultaneous rotation invariance removes the quotient.
\end{proof}

\subsection{The ground-state transform}

Define positive functions on the open chamber by
\begin{equation}\label{sp:ground-state}
 J(w)=\prod_{1\le i<j\le4}\sin\frac{w_j-w_i}{2},\qquad
 h(x,y)=J(x)^{1/3}J(y)^{1/3},\qquad \lambda=\frac54.
\end{equation}
The product $J$ is unchanged by a cyclic relabeling.  For one tuple put
$s_i=\sum_{j\ne i}\cot((w_i-w_j)/2)$.  The elementary three-angle
cotangent identity gives
\[
 \sum_i s_i^2=\sum_{i\ne j}\csc^2\frac{w_i-w_j}{2}-20,
 \qquad \partial_{w_i}\log J^{1/3}=\frac{s_i}{6}.
\]
Substitution in $L$ shows that each tuple contributes $-5/8$ to
$h^{-1}Lh$.  Hence $Lh=-\lambda h$.

This ground-state calculation is the cotangent-diffusion form of the
SLE--Calogero--Sutherland structure studied in \cite{CardyCS}; we give
the estimates needed here directly.  The transformed generator and
its invariant density are
\begin{equation}\label{sp:doob-generator}
 D=\frac38\bigl(\Delta+\nabla\log m\cdot\nabla\bigr),
 \qquad m(x,y)=J(x)^{4/3}J(y)^{4/3}.
\end{equation}
Let $\mu$ be the probability measure proportional to $m$ times Lebesgue
angle measure.  For $T=T^+,T^-$, or a linear combination of them, set
$W(p)=e^{\lambda p}T(p)/h$.  Equation~\eqref{sp:pde} becomes
\begin{equation}\label{sp:transformed-pde}
 \partial_pW=DW+B_p\cdot\nabla W+V_pW,
 \qquad V_p=B_p\cdot\nabla\log h.
\end{equation}
The apparent singularities of $V_p$ at collisions are removable.
Every cotangent is multiplied by a difference of two components of
$B_p$, which vanishes on the corresponding collision face.
Periodicity handles the face $w_4=w_1+2\pi$.  Thus $V_p$ has smooth
bounds $O(e^{-p})$ up to the chamber boundary, as does $B_p$.

\begin{lemma}[Smoothing on the chamber]\label{sp:smoothing}
Fix $p_0>0$.  For either transmission function, or its sum or difference,
the transformed solution satisfies
\begin{equation}\label{sp:uniform-smoothing}
 \|W(p)\|_\infty+\operatorname{Lip}(W(p))\le C_{p_0}
 \qquad(p\ge p_0+2).
\end{equation}
The semigroup $P_t$ of $D$ is symmetric on $L^2(\mu)$ and, for every
$q>1$, maps $L^q(\mu)$ to bounded functions at every fixed positive
time.  At unit time the evolution in \eqref{sp:transformed-pde} maps
bounded functions to Lipschitz functions with a constant uniform in
the starting modulus $p\ge p_0$.
\end{lemma}

\begin{proof}
We first justify using the diffusion all the way to the chamber
boundary.  The drift in one tuple is
$\frac14(\sum_{j\ne i}\cot((w_i-w_j)/2))_{i=1}^4$.
On the lifted convex chamber, use $-\log m$, shifted to be nonnegative,
as a Lyapunov function.  If $\alpha=4/3$ is the density exponent, the
cotangent identity shows that its generator has a singular contribution
proportional to $(\alpha-\alpha^2)\sum_{i<j}\csc^2((w_i-w_j)/2)$.
This coefficient is negative, so the generator is bounded above.
More explicitly, the one-tuple contribution is
\[
 D(-\log J^{4/3})=
 \frac{10}{3}-\frac1{24}\sum_{i\ne j}
                           \csc^2\frac{w_i-w_j}{2}.
\]
Stopping on inner exhaustions proves that no collision occurs in
finite time.  Interior existence and pathwise uniqueness therefore
give a global diffusion.  A bounded added drift, such as the
time-dependent $B_p$, preserves noncollision by Girsanov's theorem on
each bounded time interval.  The divergence form of $D$ gives symmetry
and invariance of $\mu$, first on interior exhaustions and then by
noncollision.

The Feynman--Kac representation for
\eqref{sp:transformed-pde} uses $D+B\cdot\nabla$ and potential $V$,
with coefficients read at decreasing moduli from the evaluation
modulus down to the initial modulus.  Its initial data at $p_0$ obey
$|W(p_0)|\le C/h$.  To justify this representation despite the
unbounded transformed data, first apply the killed representation to
$T$ on inner chamber exhaustions.  Lemma~\ref{an:collision} says that
the lateral values of $T$ tend uniformly to zero on each compact
modulus interval.  Their contribution therefore disappears.
The $h$-transform on the exhaustion gives the stated representation;
pass to the limit using noncollision.  It suffices to do this for the
nonnegative $T^\pm$ and then use linearity.

Here are the uniform smoothing estimates.  The negative logarithm of
$m$ is convex on the lifted chamber, so the drift of $D$ is monotone
dissipative.  Couple two diffusions by the same Brownian motion and
add to the second a drift toward the first of magnitude their initial
distance divided by the allotted time, until they meet.  Dissipativity
forces meeting by that time.  Girsanov's theorem removes the extra
drift with all fixed moments bounded in terms of the time and the
initial distance.  Localization inside the chamber justifies the
coupling, and the preceding change-of-measure argument excludes
collisions for the controlled diffusion.  H\"older's inequality gives
the dimension-free form
\[
 |P_tF(z)|^q\le
 C(q,t,\operatorname{dist}(z,z'))P_t(|F|^q)(z').
\]
One can always choose lifts of $z,z'$ at bounded distance.
Integration in $z'$ against the invariant probability $\mu$ proves
the asserted $L^q$-to-supremum estimate.

A bounded drift perturbation and a bounded potential on a unit time
interval preserve initial smoothing from $L^4(\mu)$.
Indeed Cauchy--Schwarz applied to the Girsanov density bounds the
perturbed expectation of $|f|$ by $C(P_1|f|^2)^{1/2}$;
the unperturbed $L^2$-to-supremum bound for $|f|^2$ then gives
$C\|f\|_{L^4(\mu)}$.  There is exactly this integrability here:
$h^{-1}\in L^4(\mu)$, since $h^{-4}m=1$.
After the first unit interval $W$ is bounded.  Its bound persists
for all subsequent times because
$\int_{p_0}^\infty\|V_p\|_\infty\,dp<\infty$.

The same coupling proves a Lipschitz estimate for bounded data on a
unit interval.  In the presence of $B$, a control of size at most a
fixed constant times the initial distance forces meeting; before
meeting the distance stays bounded by the same multiple of the
initial distance.  The second moment estimate for the Girsanov
exponential bounds its total-variation cost by that distance.
The two Feynman--Kac weights differ by at most the same order, since
$V$ has a uniform Lipschitz bound.  Thus the unit-time Lipschitz norm
is bounded by a constant times the initial supremum norm.  Apply this
on the last unit interval to obtain
\eqref{sp:uniform-smoothing}.
\end{proof}

\subsection{The two symmetry gaps}

The role of rotation averaging can now be seen explicitly.  For a
tuple with real ordered lift $w$, write
\[
 \bar w=\frac14\sum_iw_i,\qquad
 r=w-\bar w\mathbf1,\qquad \theta=\bar w\pmod{2\pi}.
\]
The relative coordinate lies in the bounded convex alcove
\[
 \mathcal A=\{r\in\R^4:\textstyle\sum_ir_i=0,
                         \ r_1<r_2<r_3<r_4<r_1+2\pi\}.
\]
This identifies $\mathcal C$ with a product of a circle and
$\mathcal A$.  Its invariant law is a product as well, since $J$
depends only on the relative angles.  The center part of $D$ is
$(3/32)\partial_\theta^2$ for each tuple.

\begin{lemma}[Spectral gaps]\label{sp:gaps}
The relative-angle diffusion for one tuple has Poincar\'e gap at least
$1/2$.  Consider the subspace of $L^2(\mu)$ consisting of functions
invariant under simultaneous rotation of the two tuples and changing
sign under one cyclic relabeling of either tuple.  On this subspace
$-D$ has gap at least $3/4$.  On its subspace with both center
frequencies zero the gap is at least $1$.
\end{lemma}

\begin{proof}
For a mean-zero vector $v\in\R^4$ the Hessian of the negative
logarithm of the one-tuple density satisfies
\[
 \nabla^2(-\log J^{4/3})[v,v]
 =\frac13\sum_{i<j}\csc^2\frac{r_i-r_j}{2}(v_i-v_j)^2
 \ge\frac43|v|^2.
\]
The Brascamp--Lieb inequality \cite{BL,KM}, applied on the convex alcove,
therefore gives
\[
 \operatorname{Var}(F)\le\frac34\int|\nabla_rF|^2\,d\mu_r.
\]
The Dirichlet form has factor $3/8$, yielding the gap $1/2$.
One may first use inner convex exhaustions of the alcove and then
pass to the limit.

For the two tuples, simultaneous rotation invariance allows only
center frequencies $(n,-n)$, whose energy cost is $3n^2/16$.
A cyclic relabeling sends $\theta$ to $\theta+\pi/2$ and acts on
$\mathcal A$ by a measure-preserving affine isometry.  If a function
changes sign under this relabeling, its relative-angle mean must
vanish unless $e^{in\pi/2}=-1$, that is, unless $n\equiv2\pmod4$.
The same assertion holds separately for the other tuple.  Therefore
frequencies not congruent to $2$ modulo $4$ pay at least $1/2+1/2=1$
in relative energy.  The remaining frequencies have $|n|\ge2$ and
pay at least $3/4$ in center energy.  At $(n,-n)=(0,0)$ the two
relative means vanish, giving the gap $1$.

These inequalities give $L^2$ semigroup decay through the symmetric
Dirichlet form
\[
 \mathcal E(F,F)=\frac38\int|\nabla F|^2\,d\mu.
\]
There is no hidden absorbing boundary: the constant function is in
the form domain.  Cutoffs approaching the collision faces have
vanishing energy because the density vanishes with power $4/3>1$.
\end{proof}

Let $W_-=e^{\lambda p}(T^+-T^-)/h$, and define
\[
 (\Pi F)(x,y)=\frac1{2\pi}\int_0^{2\pi}
                           F(x,y+\phi\mathbf1)\,d\phi
\]
on simultaneous-rotation-invariant functions.  Equivalently, $\Pi$
is projection to center frequencies $(0,0)$.  A one-step cyclic
relabeling switches $T^+$ and $T^-$, while $h$ is unchanged.
Thus $W_-$ has exactly the symmetries of Lemma~\ref{sp:gaps}.

\begin{proposition}[Uniform transmission asymptotics]
\label{sp:transmission-asymptotics}
There are constants $K\ge0$, $\sigma>0$, and $C<\infty$ such that,
uniformly over the open angle chamber and all sufficiently large $p$,
\begin{align}
 \|W_-(p)\|_\infty&\le Ce^{-3p/4},
 &\|\Pi W_-(p)\|_\infty&\le C(1+p)e^{-p},
 \label{sp:odd-decay}\\
 T^\pm(p,x,y)&=e^{-\lambda p}
                   \bigl(Kh(x,y)+O(e^{-\sigma p})\bigr).
 \label{sp:rank-one-round}
\end{align}
Moreover,
\begin{equation}\label{sp:odd-regularity}
 \operatorname{Lip}(W_-(p))\le Ce^{-3p/4},\qquad
 \|W_-(p+\ell)-W_-(p)\|_\infty
        \le Ce^{-3p/4}\sqrt\ell\quad(0\le\ell\le1).
\end{equation}
The function $h$ is uniformly $1/3$-H\"older on the closed chamber.
\end{proposition}

\begin{proof}
For any transformed solution put
$F_p=B_p\cdot\nabla W(p)+V_pW(p)$.
Lemma~\ref{sp:smoothing} implies $\|F_p\|_\infty\le Ce^{-p}$
once $p$ is sufficiently large.  Variation of constants gives
\begin{equation}\label{sp:duhamel}
 W(p)=P_{p-p_1}W(p_1)+\int_{p_1}^pP_{p-s}F_s\,ds.
\end{equation}
This identity also follows directly by It\^o's formula with $D$,
first stopping on inner exhaustions: on the interval in question
the solution and the forcing are bounded, and collisions are not hit.

For $W_-$ the forcing respects simultaneous rotation invariance and
the two cyclic sign changes.  Apply the $3/4$ gap to
\eqref{sp:duhamel}, and then apply the unit-time $L^2$-to-supremum
smoothing to all semigroup times at least one.  For shorter times
use the supremum bound on $F_s$.  The convolution with $e^{-s}$ is
$O(e^{-3p/4})$.  After projecting by $\Pi$, which commutes with $D$,
the gap is $1$ and the convolution is $O((1+p)e^{-p})$.
This proves \eqref{sp:odd-decay}.  It does not require $B_p$ to
commute with independent rotations: the projected forcing is
estimated after projection.

For either $W^+$ or $W^-$ separately, $D$ has a positive gap off
constants.  Indeed the center gap is $3/32$ and the relative gap is
at least $1/2$.  The same argument shows exponential convergence of
$W-\int W\,d\mu$ to zero.  Also
$\partial_p\int W\,d\mu=\int F_p\,d\mu=O(e^{-p})$, so its mean
converges exponentially.  The first estimate of
\eqref{sp:odd-decay} makes the two limiting constants equal; call
their common value $K$.  Multiplication by $e^{-\lambda p}h$
proves \eqref{sp:rank-one-round}.

The unit-interval Lipschitz estimate in Lemma~\ref{sp:smoothing},
applied to the homogeneous perturbed evolution for $W_-$, improves
its Lipschitz bound to $Ce^{-3p/4}$.  To obtain time regularity,
the same diffusion has expected angle displacement at most
$C\sqrt\ell$ in time $\ell\le1$, uniformly in its initial point.
Here is a bound that remains valid near a collision.  Apply It\^o's
formula to squared distance from the initial lift.  Each pair in
the scalar product with the singular drift contributes a constant
multiple of
\[
 (v-v_0)\cot(v/2),\qquad v,v_0\in(0,2\pi),
\]
which is bounded above uniformly: use $v-v_0\le v$ for $v\le\pi$
and $v-v_0\ge v-2\pi$ for $v\ge\pi$.
The bounded additional drift and the Brownian quadratic variation
then give $\E|Z_\ell-Z_0|^2\le C\ell$, by localization and
Gronwall's inequality.  The Feynman--Kac representation over time
$\ell$, the spatial Lipschitz bound, and the bounded potential give
the time estimate in \eqref{sp:odd-regularity}.
Finally $J$ is a bounded Lipschitz nonnegative function on the closed
chamber, and taking a cube root is $1/3$-H\"older.  The assertion
for $h$ follows by taking products.
\end{proof}
